\documentclass[11pt]{article}

\usepackage[a4paper,margin=1in]{geometry}
\usepackage{amsmath,amssymb}
\usepackage{hyperref}

\hypersetup{
    colorlinks=true,
    linkcolor=blue,
    citecolor=blue,
    urlcolor=blue
}

\title{Dirac Measure}
\author{}
\date{}

\begin{document}

\maketitle

\section*{Definition}

Let $(X,\mathcal A)$ be a measurable space and let $x\in X$. The Dirac measure
at $x$, denoted by $\delta_x$, is the measure defined by
\[
\delta_x(A)
=
\begin{cases}
1, & x\in A,\\
0, & x\notin A,
\end{cases}
\qquad A\in\mathcal A.
\]

\section*{Interpretation}

The Dirac measure places all its mass at the single point $x$. If $f:X\to\mathbb R$
is measurable, then
\[
\int_X f\,d\delta_x = f(x),
\]
whenever the integral is defined.

\section*{In Euclidean Space}

On $\mathbb R^d$ with its Borel $\sigma$-algebra, $\delta_x$ is a probability
measure concentrated at the point $x\in\mathbb R^d$.

\end{document}
